\begin{frame}{수렴이 얼마나 빠른가 --- 지수적으로}
  \begin{columns}[T]
    \begin{column}{0.55\textwidth}
      \begin{itemize}
        \item 매 반복의 변화량(전체 노드 중 최대 $|\Delta PR|$)은
          \textbf{지수적으로} 줄어든다.
        \item 카라테 클럽: $d{=}0.85$ 일 때 대략 \textbf{45회}로
          $10^{-9}$ 이내 수렴. $d{=}0.5$ 이면 \textbf{19회}.
        \item $d$가 클수록 서퍼가 더 ``\textbf{장기 기억}''을 하려
          하므로 수렴이 \textbf{느려진다}.
        \item 작~중형 그래프: ``\textbf{100회 반복}''이면 충분.
        \item 수십억 페이지 규모의 웹: 구글이 \textbf{분산 계산 ·
          희소화} 같은 별도 공학이 필요했던 이유.
      \end{itemize}
    \end{column}
    \begin{column}{0.42\textwidth}
      \centering
      \includegraphics[width=\linewidth]{ch14_3_pagerank_convergence.png}
      {\footnotesize 매 반복의 최대 변화량 max $|\Delta PR|$ 을
      로그 스케어로. $d{=}0.85$ (파랑)가 $d{=}0.5$ (주황)보다
      더 오래 걸리며, 두 곡선 모두 직선(= 지수적 감소)을 따라
      $10^{-9}$ 아래로.}
    \end{column}
  \end{columns}
\end{frame}
